\documentclass{article}

\usepackage{neurips_2026}

\usepackage[utf8]{inputenc}
\usepackage[T1]{fontenc}
\usepackage{hyperref}
\usepackage{url}
\usepackage{booktabs}
\usepackage{graphicx}
\usepackage{amsfonts}
\usepackage{nicefrac}
\usepackage{microtype}
\usepackage{xcolor}
\usepackage{enumitem}
\usepackage{fancyvrb}
\usepackage{listings}
\lstdefinestyle{kisspy}{
  language=Python,
  basicstyle=\ttfamily\scriptsize,
  keywordstyle=\color[rgb]{0.0,0.0,0.8}\bfseries,
  commentstyle=\color[rgb]{0.25,0.5,0.35}\itshape,
  stringstyle=\color[rgb]{0.73,0.13,0.13},
  showstringspaces=false,
  frame=single,
  framesep=2mm,
  numbers=left,
  numberstyle=\tiny\color{gray},
  numbersep=5pt,
  breaklines=true,
  tabsize=4,
  captionpos=b,
}
\lstnewenvironment{prompt}{%
  \lstset{language={},basicstyle=\ttfamily\scriptsize,frame=single,framesep=2mm,breaklines=true,columns=fullflexible}%
}{}

\title{KS Assistant: A Simple General-Purpose AI Agent for Software Engineering}
\author{}

\begin{document}

\maketitle

\begin{abstract}
Large language models can generate code and call tools with remarkable fluency, yet deploying them as practical software engineering assistants still exposes stubborn gaps: finite context windows, a single mistake that can derail entire sessions, agents that get stuck in dead ends, AI slop, and generated changes that are difficult to review or revert.

We present \textbf{KS Assistant}, a general-purpose assistant and integrated development environment (IDE) built on top of the \textbf{KS Agent Framework}, a simple AI agent framework of roughly 1,900~lines of code.  The framework addresses these gaps through a robust system prompt and a five-layer agent hierarchy in which each layer adds exactly one concern: budget-tracked ReAct execution, automatic continuation across sub-sessions via summarization, coding and browser tools with parallel sub-agents, persistent multi-turn chat with history recall, and git worktree isolation so every task runs on its own branch.  Both KS Assistant and the KS Agent Framework are grounded in disciplined software engineering practice; these principles are encoded directly into the agent's system prompt, enabling KS Assistant to write code that is simple, elegant, maintainable, and bug-free.

We implemented KS Assistant as a free, open-source Visual Studio Code extension that runs locally, is effective for long-horizon tasks, and supports browser automation, multimodal input, and Docker containers.  We deliberately prioritize output quality over speed: giving a frontier model adequate time to validate its own output---running linters, type checkers, and tests---dramatically reduces the low-quality code that plagues faster but less thorough agents.  The entire system was built using itself over four months, providing a continuous stress test in which any bug was patched immediately after its manifestation.  On Terminal-Bench~2.0, KS Assistant achieves a 62.2\% overall pass rate with Claude Opus~4.6~\citep{anthropic2026opus46}, comparing favorably to Claude~Code (58\%) and Cursor Composer~2 (61.7\%).  These results are achieved without benchmark-specific prompt tuning or model fine-tuning.
\end{abstract}

%----------------------------------------------------------------------
\section{Introduction}
\label{sec:intro}
%----------------------------------------------------------------------

Modern large language models (LLMs), such as Claude Opus 4.7~\citep{anthropic2026opus47}, GPT 5.5~\citep{openai2026gpt55}, and Gemini 3.1~\citep{google2026gemini31}, have demonstrated a remarkable ability to generate code, reason about software architecture, and use developer tools~\citep{chen2021evaluating,roziere2023code}.  A growing body of work has explored how to harness these capabilities for autonomous software engineering, from single-session agents that resolve GitHub issues~\citep{yang2024sweagent,wang2024openhands} to industrial products marketed as AI software and general assistants~\citep{copilot2021,cursor2024,devin2024,claudecode2025,codex2025,openclaw2025}.  Yet using an LLM as a practical software engineering assistant still exposes several stubborn gaps: context windows are finite, a single mistake can derail an entire session, agents get stuck in dead ends, models generate AI slop, and generated changes are difficult to review or revert once applied to a live codebase.

We propose the \emph{KS Agent Framework}, a simple AI agent framework containing around 1,900 lines of code for the core agent implementation.  The name ``KS'' reflects the \emph{Keep Simple} design philosophy: each layer is small, each concern is isolated, and the overall system avoids unnecessary abstraction.  We address the above-mentioned gaps through a robust system prompt (Section~\ref{sec:system-prompt} and Appendix~\ref{app:full-prompt}) and a five-layer agent hierarchy in which each layer solves exactly one concern:

\begin{enumerate}[nosep]
  \item \textbf{KS Agent} --- budget-tracked ReAct~\citep{yao2023react} loop with native function calling.
  \item \textbf{Relentless Agent} --- automatic summarization and continuation across sub-sessions.
  \item \textbf{Tool Agent} --- coding tools, browser automation, and parallel sub-agent execution.
  \item \textbf{Chat Agent} --- persistent multi-turn chat sessions with history recall.
  \item \textbf{Worktree Agent} --- git worktree isolation so every task runs on its own branch.
\end{enumerate}

To demonstrate the power of the KS Agent Framework, we implemented KS Assistant as a Visual Studio Code extension that runs locally.  It has full browser support (using open-source Chromium and Playwright), multimodal support, Docker container support, and a mobile/web app.  KS Assistant is completely free and open-source; all one needs is a model API key from a major LLM provider.  The open-source repository link is withheld to respect double-blind review.

KS Assistant has been built using itself.  The entire codebase---the KS Agent framework, the agent layers, the VS~Code extension, and the system prompt---was developed by KS Assistant operating on its own repository.  This self-hosting discipline provides a continuous-integration-style stress test: if the agent introduces a bug that impairs its ability to function, the developers immediately ask the agent to fix it by analyzing the trajectory and code.  The simplicity of the framework enabled the agent to implement features confidently without bugs and AI slop.  The five core agent classes are remarkably compact: the KS Agent comprises 413~lines, the Relentless Agent 321~lines, the Tool Agent 322~lines, the Chat Agent 125~lines, and the Worktree Agent 704~lines---a total of roughly 1,885~lines of code (excluding empty lines and comments).

We deliberately prioritize output quality over speed.  Using a weaker or cheaper model often forces the developer to discard the agent's work and retry, ultimately increasing the total cost.  Conversely, giving a frontier model adequate time to validate its own output---running linters, type checkers, and tests before declaring success---dramatically reduces slop.  We expect token costs and inference latencies to continue to fall~\citep{gao2025trae}, making this quality-first posture increasingly practical.

We evaluate on Terminal-Bench~2.0 and achieve a 62.2\% overall pass rate using Claude Opus~4.6~\citep{anthropic2026opus46}, comparing favorably to Claude~Code (58\%) and Cursor Composer~2 (61.7\%)~\citep{cursor2026composer2} on the same benchmark (Section~\ref{sec:evaluation}).  These results are achieved without benchmark-specific prompt tuning or model fine-tuning.  We show that a simple agent framework, without sophisticated agent technologies such as trajectory compaction and asynchronous multi-agent orchestration, is sufficient to build a competitive software engineering assistant.  By building KS Assistant using itself and matching or beating both the Cursor and Claude Code agents, we demonstrate that time-tested software engineering techniques and principles are invaluable for building reliable, practical agent systems.

\paragraph{Outline.}
Section~\ref{sec:architecture} presents the five-layer agent architecture.  Section~\ref{sec:evaluation} reports evaluation results on Terminal-Bench~2.0.  Section~\ref{sec:system-prompt} describes the two most distinctive aspects of the system prompt: the execution mindset and the web research protocol (the full prompt is detailed in Appendix~\ref{app:full-prompt}).  Section~\ref{sec:vscode-features} covers unique features of the VS~Code extension.  Section~\ref{sec:related} discusses related work, and Section~\ref{sec:conclusion} concludes.

%----------------------------------------------------------------------
\section{Agent Architecture}
\label{sec:architecture}
%----------------------------------------------------------------------

The KS Agent Framework was originally built to rapidly prototype and experiment with prompt optimization techniques~\citep{agrawal2025gepa} and evolutionary algorithms for algorithmic optimization~\citep{novikov2025alphaevolve,openevolve2025}.  The emphasis on simplicity enabled coding agents to write simple, bug-free code.  No prompt optimization techniques were used when creating the system prompt for KS Assistant; it was hand-tuned based on long-term experience with the agent and its behavior.  The framework uses five agent layers combining composition and inheritance.  Each layer delegates upward for concerns it does not own.

%......................................................................
\subsection{KS Agent}
\label{sec:ks-agent}
%......................................................................

The KS Agent is the innermost execution unit implementing a standard ReAct loop~\citep{yao2023react}.

\begin{lstlisting}[style=kisspy,numbers=none,caption={A complete KS agent with a single tool.},label={lst:simple-agent}]
from ks.core.ks_agent import KSAgent

def calculate(expression: str) -> str:
    """Evaluate a math expression."""
    return str(eval(expression))

agent = KSAgent(name="Math Buddy")
result = agent.run(
    model_name="gemini-2.5-flash",
    prompt_template="Calculate: {question}",
    arguments={"question": "What is 15% of 847?"},
    tools=[calculate]
)
print(result)  # 127.05
\end{lstlisting}

\textbf{Native function calling.}  Tools are registered as ordinary Python callables.  The agent builds an OpenAI-compatible tool schema once at setup time and caches it.  A special \texttt{finish} tool signals task completion and returns the result.

\textbf{Step, token, and budget tracking.}  At every step, the agent extracts input and output token counts from the API response, computes dollar cost using a per-model pricing table, and updates both local and global budget counters protected by a class-level lock.  Three limits are checked before each step: per-agent budget, global budget, and maximum step count.

\textbf{Error resilience.}  The agent retries transient API errors (rate limits, server errors) up to a configurable threshold.  Non-retryable errors (authentication failures, permission denials) are raised immediately.

\textbf{Non-agentic mode.}  When tools are not needed, the agent runs a single generation without the ReAct loop, useful for summarization sub-tasks.

Listing~\ref{lst:simple-agent} shows a complete, working agent in under ten lines of code.  The developer defines an ordinary Python function (\texttt{calculate}), instantiates a \texttt{KSAgent}, and calls its \texttt{run} method with a model name, a prompt template, template arguments, and a list of tools.  The framework automatically handles tool-schema generation, the ReAct loop, and budget tracking.

The KS Agent is stateless across runs: each call resets the conversation, token counters, and tool registry, making it safe to reuse a single instance for multiple sequential tasks.

%......................................................................
\subsection{Relentless Agent}
\label{sec:relentless-agent}
%......................................................................

The Relentless Agent wraps a KS Agent in a continuation loop.  Its core contribution is executing tasks that exceed a single context window by breaking them into sub-sessions.

Rather than investing in context-compaction techniques, we adopt a simple continuation protocol: when a sub-session exhausts its context window or step budget, the agent produces a \emph{structured summary} of every action taken so far---chronologically ordered, with explanations and relevant code snippets---and a fresh sub-session resumes from that summary.  This approach is related to Reflexion~\citep{shinn2023reflexion}, which feeds verbal self-critiques back into subsequent trials, but uses a Reflexion-like technique to \emph{continue} a task rather than retry it.  We found that a na\"ive instruction to ``summarize the current context'' produced poor continuations; requiring a \emph{step-by-step chronological account with code snippets} dramatically improved coherence across sub-sessions.  A potential limitation is that summaries may grow unwieldy for multi-day tasks; in practice, we have not encountered this problem even for tasks spanning several hours, but a thorough evaluation of summary scaling remains future work.

\textbf{Continuation protocol.}  The \texttt{finish} tool accepts three fields: a success flag, a continue flag, and a summary.  When \texttt{is\_continue=True}, the Relentless Agent starts a new sub-session with a fresh context window.  The prompt for the new session includes a chronologically ordered list of all prior attempt summaries and instructs the agent not to redo completed work:

\begin{prompt}
# Task Progress (Continuation {continuation_number})

{progress_text}

# Continue
- Complete the rest of the task.
- **DON'T** redo completed work.
- If you have been retrying the same approach without
  progress, step back and rethink the strategy from
  scratch.
\end{prompt}

\textbf{Forced continuation on failure.}  If a sub-session raises an exception (e.g., the step limit is hit before calling finish), the Relentless Agent saves the full trajectory to a temporary file, spawns a separate summarizer agent to produce a concise summary, and uses that summary as the progress text for the next sub-session.  This ensures that even crashed sessions contribute useful context.  The summarizer receives the following prompt:

\begin{prompt}
# Summarizer

The trajectory of the agent is stored in the file:
{trajectory_file}

# Instructions
- Read the trajectory file and analyze it.  The trajectory
  file could be large.
- Return a precise chronologically-ordered list of things
  the agent did, with the reason for doing that, along
  with relevant code snippets
\end{prompt}

\textbf{Pre-emptive continuation.}  To force the agent to self-continue before exhausting its step budget, the system prompt is augmented with a step-threshold instruction:

\begin{prompt}
# MOST IMPORTANT INSTRUCTIONS
- **At step {step_threshold}: you MUST call
  finish(success=False, is_continue=True,
  summary="precise chronologically-ordered list of things
  the agent did with the reason for doing that along with
  relevant code snippets")** or if the task is not complete
  and you are at risk of running out of steps or context
  length.
\end{prompt}

\noindent This instruction is injected near the end of the budget window.  Without it, the agent exhibits a ``rush to finish'' behavior when steps are running low---skipping verification, making hasty edits, and calling \texttt{finish} with an incomplete result.  The explicit step threshold redirects that urgency into the continuation protocol, ensuring that a clean handoff to a new sub-session produces better results than a frantic attempt to squeeze everything into the remaining steps.

%......................................................................
\subsection{Tool Agent}
\label{sec:tool-agent}
%......................................................................

The Tool Agent adds the tools that make the system useful for software development and general-purpose automation.

\textbf{Coding tools.}  Four core tools: a shell command executor with streaming output, a file reader, a precise string-based file editor, and a file writer.  The shell executor supports configurable timeouts, streams output in real time, and respects a stop event for user cancellation.

\textbf{Browser automation.}  A web-use tool provides programmatic browser control: navigating to URLs, reading page accessibility trees, clicking elements, typing text, pressing keys, scrolling, and taking screenshots.  It uses the open-source Chromium browser via the Playwright library.

\textbf{Parallel sub-agents.}  An optional parallel execution tool spawns independent Tool Agent instances in a thread pool.  Each sub-agent gets its own LLM context and tool set, useful for embarrassingly parallel tasks such as summarizing multiple files or researching independent topics.  Results are collected and returned in input order.  This tool is not activated by default because the IDE cannot stream multiple agent outputs coherently in the chat window.

\textbf{User interaction.}  An ask-user-question tool pauses execution and requests clarification from the user.

\textbf{Docker isolation.}  When a Docker image is specified, coding tools are replaced with Docker-aware variants that execute commands inside a container.

%......................................................................
\subsection{Chat Agent}
\label{sec:chat-agent}
%......................................................................

The Chat Agent adds multi-turn conversation persistence.

\textbf{Chat sessions.}  Each task is assigned to a chat session identified by a stable chat ID.  Tasks and results are persisted to a local database.  New tasks within the same session include prior tasks and results as numbered context entries, allowing the LLM to reference earlier work.

\textbf{Bounded chat context.}  The agent caps in-context history entries at $K=10$.  When the cap is exceeded, it preserves the first two entries (which establish the user's intent) and the most recent entries, dropping the middle entries least likely to be referenced.

\textbf{Session management.}  Three operations are supported: starting a new chat, resuming by task description, and resuming by explicit chat ID.

\textbf{Frequent task tracking.}  Each time a task is executed, the agent records the task description in a frequency table.  The IDE sidebar surfaces the most frequent tasks so users can re-issue them with one click, turning recurring requests (``run the test suite,'' ``regenerate the changelog'') into a click-to-replay experience.

\textbf{Metadata persistence.}  After each task, the agent records metadata including the model used, working directory, software version, token counts, cost, and whether the task used parallel execution or worktree isolation.  This audit trail supports cost analysis and debugging.

%......................................................................
\subsection{Worktree Agent}
\label{sec:worktree-agent}
%......................................................................

The Worktree Agent is the outermost layer.  Its defining feature is git-worktree isolation.

\textbf{Branch-per-task.}  When a task starts, the agent creates a new git branch and corresponding worktree directory.  All agent modifications happen inside the worktree; the user's main working tree remains untouched.

\textbf{Dirty-state preservation.}  Uncommitted changes in the main working tree are copied into the worktree with a baseline commit.  During merge, cherry-pick from the baseline replays only the agent's changes.

\textbf{Concurrency safety.}  A per-repository file lock serializes checkout, stash, merge, and pop operations.  Thread-local storage isolates per-task state.

\textbf{Crash recovery.}  All worktree state is stored in git itself (branch names, git config entries) rather than sidecar files.  On process restart, the agent queries git and reconstructs instance attributes, enabling seamless recovery.

\textbf{Graceful fallback.}  If the working directory is not inside a git repository, has no commits, or has a detached HEAD, the agent falls back to direct execution without worktree isolation.

%----------------------------------------------------------------------
\section{Evaluation on Terminal-Bench~2.0}
\label{sec:evaluation}
%----------------------------------------------------------------------

We evaluate on Terminal-Bench~2.0,\footnote{\url{https://www.tbench.ai/}} a benchmark comprising 89 diverse terminal-based programming tasks, ranging from building legacy compilers and configuring servers to solving cryptanalysis challenges and training machine-learning models.  Each task runs in an isolated Docker container; a separate verifier automatically judges the result.  We use the Harbor\footnote{\url{https://github.com/harbor-framework/harbor}} framework to orchestrate execution and Claude Opus~4.6~\citep{anthropic2026opus46} as the underlying LLM.  We do not modify the general system prompt or inject benchmark-specific instructions.  Evaluation was carried out on a 2025 MacBook Air 15'' with an M4 processor and 24\,GB RAM.

\subsection{Setup}

We run 5 independent trials per task.  The agent is a thin Harbor adapter that installs and invokes the KS Assistant CLI inside each container.  We hard-skip 9 tasks verified to be infeasible for Opus~4.6 across 6+ prior attempts (e.g., CompCert compilation, Windows~3.11 GUI installation, video OCR) to save time and token cost.  Skipped tasks still count as failures.

\subsection{Aggregate Results}

Table~\ref{tab:tb2-aggregate} summarizes the aggregate statistics.

\begin{table}[h]
\centering
\caption{Terminal-Bench~2.0 aggregate results (89 tasks, 5 trials each, Claude Opus~4.6).}
\label{tab:tb2-aggregate}
\begin{tabular}{lr}
\toprule
\textbf{Metric} & \textbf{Value} \\
\midrule
Total tasks              & 89 \\
Overall pass rate        & 62.2\% (277/445) \\
pass@any (at least 1/5 passes) & 78.7\% (70/89) \\
pass@all (all 5 pass)    & 43.8\% (39/89) \\
Always-fail tasks        & 19 \\
Always-pass tasks        & 39 \\
Mixed-result tasks       & 31 \\
Median cost per trial    & \$0.45 \\
Mean cost per trial      & \$0.90 \\
Median duration per trial & 202\,s \\
Mean duration per trial  & 446\,s \\
\bottomrule
\end{tabular}
\end{table}

The 62.2\% overall pass rate compares favorably to other agents using the same underlying model: at the time of writing, Claude~Code (also Opus~4.6) scores approximately 58\% on the Terminal-Bench~2.0 leaderboard, and Cursor's Composer~2---a custom fine-tuned model trained with large-scale reinforcement learning~\citep{cursor2026composer2}---achieves 61.7\%.  Our result suggests that the layered architecture and the structured system prompt contribute meaningfully beyond what the base model provides.

\subsection{Task-Level Breakdown}

\textbf{Consistently solved tasks (39 of 89).} These include cryptanalysis (FEAL differential), game-playing (chess best move), git operations (leak recovery), server configuration (gRPC key-value store, PyPI server, NGINX logging), data processing (resharding), formal verification (Coq~\texttt{plus\_comm}), ML inference (HuggingFace model serving, LLM batching scheduler), and system emulation (QEMU startup).  The breadth of this set demonstrates that the agent's generality is not limited to a single domain.

\textbf{Consistently failed tasks (19 of 89).} The failures cluster into three categories: (1)~tasks requiring graphical or multimedia capabilities unavailable in the container (video processing, Windows~3.11 GUI installation, MTEB leaderboard scraping), (2)~tasks demanding very long or resource-intensive builds that exceed time or memory limits (CompCert compilation, Doom for MIPS, Caffe CIFAR-10, training fastText on Yelp data), and (3)~tasks with niche domain-specific requirements that the model struggles to satisfy (DNA insertion, OCaml GC patching, polyglot C/Python binaries, protein assembly, cell segmentation).

\textbf{Mixed-result tasks (31 of 89).} Tasks such as \texttt{write-compressor} (3/5), \texttt{crack-7z-hash} (4/5), and \texttt{feal-linear-cryptanalysis} (4/5) succeed in most trials but occasionally fail due to non-determinism in the model's reasoning or timing-sensitive environment interactions.  Conversely, \texttt{cancel-async-tasks} (1/5) and \texttt{dna-assembly} (1/5) succeed rarely, suggesting they are at the boundary of the model's capability.

\textbf{Leaderboard context.}  KS Assistant does not score as high as some coding agents on the Terminal-Bench~2.0 leaderboard.  However, we note that recent analysis has found widespread cheating on popular agent benchmarks, including Terminal-Bench~2.0: the top three submissions commit harness-level cheating (e.g., leaking verifier code or answer keys into the agent's environment), and task-level cheating (e.g., Googling answers, mining git history, hardcoding test outputs) affects 28+ submissions across 9 benchmarks~\citep{stein2026cheatingblog,stein2026detecting}.  Separately, an automated benchmark audit found 45 confirmed hacking solutions across 13 widely used benchmarks exhibiting process-isolation failures, answer leakage, and weak test assertions that allow perfect scores without solving a single problem~\citep{wang2026trustworthyblog}.

%----------------------------------------------------------------------
\section{The System Prompt}
\label{sec:system-prompt}
%----------------------------------------------------------------------

The system prompt is a structured document that governs the agent's behavior across all tasks.  It is not a generic instruction to ``be helpful'' but rather a precise specification of engineering discipline.  We present the two most distinctive aspects here; the complete prompt is detailed in Appendix~\ref{app:full-prompt}.

%......................................................................
\subsection{Execution Mindset}
\label{sec:execution-mindset}
%......................................................................

The prompt opens with two directives that set the tone for the entire session:

\begin{prompt}
# FOCUS ON THE GIVEN TASK. ITS COMPLETION IS YOUR SOLE GOAL.
# BE RELENTLESS. BE CALM. BE RIGOROUS. BE ACCURATE. CHECK FACTS. NO AI SLOP.
\end{prompt}

\noindent Each directive addresses a specific failure mode observed in LLM-based agents:

\textbf{``FOCUS ON THE GIVEN TASK.\ ITS COMPLETION IS YOUR SOLE GOAL.''}  LLMs have a tendency to drift: they comment on the difficulty of a problem, explore tangential concerns, or ask clarifying questions that the user has already answered.  This directive anchors the model on the task at hand and discourages meta-commentary that consumes tokens without making progress.

\textbf{``BE RELENTLESS.''}  When an agent encounters an error---a failing test, a type violation, a command that returns unexpected output---the default LLM behavior is to apologize, summarize the failure, and ask the user what to do next.  This directive instructs the model to treat errors as obstacles to overcome, not as reasons to stop.

\textbf{``BE CALM.''}  The complement of relentlessness.  An overly aggressive agent may thrash between approaches without deliberation.  This directive encourages the model to analyze errors methodically before reacting.

\textbf{``BE RIGOROUS.''}  This instructs the model to follow the verification and testing disciplines encoded later in the prompt, rather than taking shortcuts when a solution \emph{looks} right.

\textbf{``BE ACCURATE.''}  LLMs frequently generate plausible-but-wrong code, file paths, or command-line flags.  This directive raises the model's threshold for asserting facts, encouraging it to verify claims against the actual file system or documentation rather than relying on parametric memory.

\textbf{``CHECK FACTS.''}  A more specific version of ``be accurate,'' targeting information the agent collects via web tools.

\textbf{``NO AI SLOP.''}  ``AI slop'' refers to the low-quality, generic, hedging text that LLMs produce when they lack confidence: filler phrases like ``certainly!'', unnecessary caveats, and boilerplate explanations.  This directive encourages the model to be concise and substantive.

\medskip
\noindent These two sentences occupy only two lines, yet they establish a behavioral contract that permeates every subsequent interaction.  We observe empirically that removing any single directive degrades the agent's behavior along the corresponding dimension (e.g., removing ``BE RELENTLESS'' causes the agent to give up after the first error more frequently).

%......................................................................
\subsection{Web Research Protocol}
\label{sec:web-research}
%......................................................................

When the agent needs external knowledge, the prompt prescribes a structured research workflow rather than allowing ad-hoc browsing:

\begin{prompt}
## Web Research (MANDATORY)

- **Visit >=30 websites every search. Hard requirement---don't stop before 30 or rationalize fewer.**
- Procedure:
  1. Create PWD/tmp/information-{unique_id}.md: `# Web Research --- Websites visited: 0/30`
  1. Per site, append: `## [N/30] URL` + extracted info. Update header counter each visit.
  1. **Don't proceed until counter >=30.**
  1. If results dry up, try different queries, synonyms, official docs, GitHub repos/issues, Stack Overflow, blogs, Reddit, papers, API refs.
  1. After 30, review and think deeply.
- Ask user for login help when needed.
\end{prompt}

\noindent The rationale is a two-phase separation between \emph{collection} and \emph{synthesis}.  LLMs tend to anchor on the first few results they encounter, biasing their solutions toward a narrow slice of the design space.  By forcing the agent to accumulate a broad set of information into a file \emph{before} reasoning about it, the protocol counteracts anchoring bias and encourages the model to consider diverse approaches.

The ``visit $\geq$30 websites'' threshold is deliberately aggressive: it ensures the agent does not shortcut the research phase after two or three hits, and the resulting information file serves as an auditable artifact of what the agent considered.  The structured procedure with a counter header and per-site entries addresses an empirically observed failure mode in which the model claims to have ``visited many sites'' after only a handful of fetches; a concrete counter forces the model to verify the actual number visited before declaring the collection phase complete.  The instruction to try ``different queries, synonyms, official docs'' when results dry up prevents the agent from giving up prematurely on a narrow set of search terms.

This two-phase collect-then-synthesize discipline is also applied to local file browsing (Appendix~\ref{app:file-browsing}): the agent writes a structured summary of each file's relevant information into a temporary markdown file, externalizing its understanding into a compact artifact that persists across context boundaries.  The instruction to collect ``without overthinking'' is deliberate: during the collection phase, the agent extracts and records facts rather than analyzes.  Analysis happens in the second phase, when the agent reads its own summary and reasons about the collected information as a whole.

%----------------------------------------------------------------------
\section{VS Code Extension: Unique Features}
\label{sec:vscode-features}
%----------------------------------------------------------------------

We release our system as a VS~Code extension and a web app.  While the underlying agent architecture already differs from existing AI coding assistants, the extension introduces several features that, to our knowledge, are not present in competing assistants---including GitHub Copilot~\citep{copilot2021}, Cursor~\citep{cursor2024}, Windsurf~\citep{windsurf2024}, Devin~\citep{devin2024}, and Aider~\citep{aider2023}.

\textbf{Cross-session self-improvement.}  The extension maintains a \texttt{USER\_PREFS.md} file that the agent reads at the start of each task and \emph{updates} at the end.  When the agent discovers a new preference or project invariant during task execution, it writes it to the file.  The agent also resolves conflicts: if a new preference contradicts an existing one, the old entry is removed.  Over time, the file accumulates a curated set of project-specific knowledge, making the agent progressively more effective without any manual configuration.  This mechanism is detailed in Appendix~\ref{app:self-improvement}.

\textbf{Real-time budget accountability.}  The extension displays real-time cost tracking in the sidebar: input tokens, output tokens, cache hits, dollar cost, and elapsed time are updated at every agent step.  Both per-task and global budget ceilings are enforced.

\textbf{Integrated browser automation.}  The extension includes a browser automation tool that allows the agent to navigate URLs, read accessibility trees, click elements, type text, press keys, scroll, and take screenshots---all controlled programmatically from within a VS~Code task.

%----------------------------------------------------------------------
\section{Related Work}
\label{sec:related}
%----------------------------------------------------------------------

\textbf{Code-specialized language models.}  Code Llama~\citep{roziere2023code} fine-tunes Llama~2 for code generation.  StarCoder~\citep{li2023starcoder} trains on permissively licensed code.  DeepSeek-Coder~\citep{guo2024deepseek} trains on a 2-trillion-token corpus.  Frontier models such as Claude Opus~4.7~\citep{anthropic2026opus47}, GPT~5.5~\citep{openai2026gpt55}, Kimi K2.5~\citep{kimiteam2026k25}, and GLM-5.1~\citep{zai2026glm51} target agentic software engineering.  Our system is model-agnostic and benefits from advances in code-specialized pre-training without architectural changes.

\textbf{Code generation agents.}  SWE-Agent~\citep{yang2024sweagent} and OpenHands~\citep{wang2024openhands} provide LLM-based agents for resolving software engineering tasks.  Agentless~\citep{xia2024agentless} shows that a two-phase localize-then-repair pipeline can achieve competitive results.  Devin~\citep{devin2024}, Claude Code~\citep{claudecode2025}, Codex~\citep{codex2025}, and Aider~\citep{aider2023} offer various approaches to autonomous coding.  Cursor released Composer~2~\citep{cursor2026composer2}, a fine-tuned coding model trained with large-scale reinforcement learning.  Our Relentless Agent layer addresses the single-session limitation common to most of these systems, while our worktree isolation provides stronger safety guarantees.

\textbf{ReAct, reasoning, and planning.}  ReAct~\citep{yao2023react} interleaves reasoning and action.  Chain-of-thought prompting~\citep{wei2022chain} elicits step-by-step reasoning.  Tree of Thoughts~\citep{yao2023tree} generalizes to search over reasoning paths.  Reflexion~\citep{shinn2023reflexion} introduces verbal reinforcement learning.  Our continuation protocol is conceptually related to Reflexion, but uses summaries to continue tasks rather than retry them.

\textbf{Multi-agent systems.}  ChatDev~\citep{qian2024chatdev} and MetaGPT~\citep{hong2024metagpt} model software development as conversations between role-playing agents.  AutoGen~\citep{wu2023autogen} provides a general multi-agent framework.  We use a single agent with broad tool access and optional parallel sub-agents, prioritizing practical utility over role-playing fidelity.

\textbf{Agent frameworks.}  LangChain~\citep{langchain2022}, DSPy~\citep{khattab2024dspy}, CrewAI~\citep{crewai2024}, smolagents~\citep{smolagents2025}, the OpenAI Agents SDK~\citep{openaiagentssdk2025}, and Google's ADK~\citep{googleadk2025} provide general-purpose agent infrastructure.  Our architecture is purpose-built for software engineering, with each layer addressing a specific practical concern.

\textbf{Benchmarks and evaluation.}  SWE-bench~\citep{jimenez2024swebench}, HumanEval~\citep{chen2021evaluating}, LiveCodeBench~\citep{jain2024livecodebench}, and SWE-bench Pro~\citep{deng2025swebenchpro} evaluate coding agents at various scales.  \citet{prathifkumar2025swebenchverified} raises concerns about benchmark contamination.  We evaluate on Terminal-Bench~2.0, which tests diverse terminal-based tasks in isolated Docker containers.

\textbf{Agentic software engineering.}  \citet{hassan2025agentic} articulate the foundational pillars of agentic software engineering.  \citet{li2025rise} surveys the landscape of autonomous coding agents.  \citet{wang2025agentic} surveys AI agentic programming techniques.  \citet{chatlatanagulchai2025readmes} studies agent context files.  \citet{mohsenimofidi2025context} investigates context engineering for AI agents.  Our system prompt and per-repository override mechanisms are instances of the context-engineering patterns these works advocate.

\textbf{Self-improvement and community prompts.}  \citet{robeyns2025selfimproving} demonstrates a self-improving coding agent.  Our self-improvement loop captures a lighter-weight form of this idea via accumulated preferences.  \citet{gao2025trae} applies test-time compute scaling to software engineering.  Get Shit Done (GSD)~\citep{gsd2025} is a meta-prompting system for coding agents that addresses context rot through phased planning documents.  Parts of our system prompt were inspired by this work.

%----------------------------------------------------------------------
\section{Conclusion}
\label{sec:conclusion}
%----------------------------------------------------------------------

We have shown that a simple layered, single-concern architecture can address the practical challenges of deploying LLM agents for real-world software development.  On Terminal-Bench~2.0, a benchmark of 89 diverse terminal-based tasks evaluated across 5 trials each, our system achieves a 62.2\% overall pass rate (277/445)---outperforming Claude~Code (58\%) and Cursor Composer~2 (61.7\%) on the same benchmark with the same underlying model (Claude Opus~4.6~\citep{anthropic2026opus46}).  These results are achieved without benchmark-specific optimizations, fine-tuning, or reinforcement learning: the entire agent is roughly 1,900 lines of straightforward Python.

We complement the architecture with a system prompt that encodes engineering discipline directly into the agent's behavior: read before writing, test before fixing, plan before executing, verify before finishing.  These are not novel insights---they are the practices of careful software engineering, translated into instructions that an LLM can follow.  The evaluation suggests that giving a frontier model the time and tools to validate its own output matters more than model-level customization.

\paragraph{Limitations.}  Our evaluation uses a single frontier model and benchmark; broader evaluation across open-weight models and benchmarks such as SWE-bench Pro is future work.  The continuation protocol's summaries may lose fidelity over many sub-sessions, and the quality-over-speed posture may not suit latency-sensitive workflows.

%----------------------------------------------------------------------
% References
%----------------------------------------------------------------------

{\small
\bibliographystyle{plainnat}
\bibliography{ks_assistant}
}

%----------------------------------------------------------------------
\appendix
%----------------------------------------------------------------------

\section{Full System Prompt Details}
\label{app:full-prompt}

This appendix presents the remaining sections of the system prompt not discussed in the main body.  The execution mindset and web research protocol are covered in Sections~\ref{sec:execution-mindset} and~\ref{sec:web-research}, respectively.

%......................................................................
\subsection{Tool Rules}
\label{app:tool-rules}
%......................................................................

Tool usage rules are explicit and mechanical:

\begin{prompt}
# Rules

- PWD = current working directory. Write() for new files; Edit() for small changes.
- Run Bash synchronously with `timeout_seconds` (default 300s). Retry with higher timeout on timeout. For >10 min commands, run in background, redirect output to file, poll periodically.
- Use go_to_url() for browser. Search the internet extensively.
- **User only sees the finish() summary. Include full details/results/outputs. Never include meta-descriptions like "Answered the user's question about X" or "Fixed the bug in Y".**
- Read large files in chunks. Temp files in PWD/tmp; clean up after.
- Use ULTRA thinking ALWAYS.
- **If running out of context/steps, don't rush---call finish(is_continue=true).**
\end{prompt}

Each rule addresses a specific failure mode:

\textbf{``Write() for new files; Edit() for small changes.''}  Without this distinction, the model may use \texttt{Write()} to overwrite an existing file with a slightly modified version, losing content it forgot to include.

\textbf{Bash timeout guidance.}  LLMs frequently launch shell commands without considering their runtime.  The instruction to use 300~seconds as the default, retry with higher timeouts, and run long-running commands in the background provides a mechanical protocol that handles common cases.

\textbf{``Use go\_to\_url() for browser.\ Search the internet extensively.''}  The first clause specifies which tool to use for browser-based interactions.  The second encourages the agent to proactively search the web when it encounters unfamiliar APIs or error messages, rather than relying on parametric memory.

\textbf{``User only sees the finish() summary.''}  This forces the model to put all user-facing information into the \texttt{finish} summary, preventing situations where the model ``tells'' the user something in an intermediate message and then assumes it was seen.  The companion clause about meta-descriptions prevents vague summaries, requiring concrete details and outputs.

\textbf{``Read large files in chunks.\ Temp files in PWD/tmp; clean up after.''}  Reading a 10,000-line file in a single tool call consumes a large fraction of the context window.  Chunked reading preserves capacity for the rest of the task.  Centralizing temporary files in a known directory makes cleanup predictable and prevents project-tree pollution.

\textbf{``Use ULTRA thinking ALWAYS.''}  This activates the model's extended reasoning mode, which is particularly valuable for complex, multi-step tasks.

\textbf{``If running out of context/steps, don't rush---call finish(is\_continue=true).''}  When the context window is nearly full, LLMs exhibit a ``rush to finish'' behavior, skipping verification steps.  This instruction redirects that urgency into the continuation protocol (Section~\ref{sec:relentless-agent}).

%......................................................................
\subsection{Pre-flight Checks}
\label{app:preflight}
%......................................................................

Before modifying any file, the agent is instructed to read it first:

\begin{prompt}
## Pre-flight Checks

- Read every file before modifying it. Read relevant sources if the task depends on existing architecture.
- If referenced files/commands/config don't exist, stop and ask or report---don't guess.
- **When fixing bugs/issues/races: write tests to confirm first, then fix.**
\end{prompt}

\textbf{``Read every file before modifying it.''}  The most common source of agent-introduced bugs is modifying a file based on an incorrect assumption about its current contents.  By requiring a fresh read immediately before any edit, the instruction ensures the model operates on the file's current state.

\textbf{``When fixing bugs/issues/races: write tests to confirm first, then fix.''}  This mandates test-first discipline: a test that reproduces the bug provides concrete verification that the fix is correct.  Writing the test forces the model to understand the bug precisely before attempting a fix.

%......................................................................
\subsection{Code Style Guidelines}
\label{app:code-style}
%......................................................................

The prompt encodes a minimalist code philosophy:

\begin{prompt}
## Code Style

- Simple, clean, readable code with minimal indirection. Organize in multiple files by functionality.
- Avoid unnecessary attributes, locals, config vars, tight coupling, and attribute redirections.
- DO NOT USE CLOSURES. No redundant abstractions or duplicate code.
- Public methods MUST have full documentation.
- Fix root causes, not symptoms. Think first: is the code simple, elegant, general, minimal?
- Don't write documentation unless the task requires it.
\end{prompt}

Each guideline addresses an anti-pattern commonly exhibited by LLM-generated code.  For example, ``DO NOT USE CLOSURES'' steers the model toward explicit data structures (plain functions with arguments, classes with attributes) that are easier to test and reason about.  ``Think first: is the code simple, elegant, general, minimal?'' is a metacognitive instruction that asks the model to pause and evaluate its plan, spending more inference-time compute on design.

%......................................................................
\subsection{Deep Work Rules}
\label{app:deep-work}
%......................................................................

\begin{prompt}
## Deep Work

- For "align"/"match"/"make consistent": read the target state before editing. Never edit from vague references.
- Use concrete values, not indirections (read Y first, then write specific values into X).
- List concrete planned changes before executing multi-part work.
- Every meaningful change needs a concrete verification method (test, grep, CLI).
\end{prompt}

These rules address failures where the model interprets instructions loosely and makes changes that are directionally correct but concretely wrong.

%......................................................................
\subsection{Planning for Complex Tasks}
\label{app:planning}
%......................................................................

\begin{prompt}
## Complex Task Planning

For 3+ files, cross-module, or architectural work:

1. List files to change and why.
1. State exact intended change per file.
1. Identify dependencies and execution order.
1. State verification method per change.

Skip for simple single-file tasks.
\end{prompt}

The complexity threshold---three or more files---avoids the overhead of planning trivial changes while ensuring comprehensive planning for tasks with cross-module dependencies.

%......................................................................
\subsection{Testing Instructions}
\label{app:testing}
%......................................................................

\begin{prompt}
## Testing

- Run lint/typecheckers; fix all errors. Achieve 100% branch coverage. Every error, including pre-existing ones, is yours---don't skip.
- NO mocks, patches, fakes, or test doubles. Write integration/e2e tests. Each test independent, verifying actual behavior.
- **Only run impacted tests after modifications.**
- To confirm races: add random sleep (<0.1s) before racing statements.
\end{prompt}

The most opinionated rule is the no-mocks requirement.  Mock tests verify that code calls certain methods in a certain order---they test the implementation, not the behavior.  A test suite built on mocks can pass with flying colors while the system is fundamentally broken.  Integration tests that exercise actual behavior provide genuine confidence that the system works.  The clause ``Every error, including pre-existing ones, is yours---don't skip'' makes the agent responsible for the entire codebase health, not just the delta it introduced.

%......................................................................
\subsection{Self-Improvement Loop}
\label{app:self-improvement}
%......................................................................

The agent maintains a preferences file that captures user invariants discovered during task execution:

\begin{prompt}
## Self-Improvement Loop

Read PWD/USER_PREFS.md at task start. Update with user preferences/invariants (no code snippets/symbols; skip one-off tasks). Remove conflicting old entries carefully and thoroughly.
\end{prompt}

This mechanism allows the agent to accumulate project knowledge over time without requiring the user to repeat themselves.  The conflict-resolution clause mandates that the agent actively resolves contradictions when updating the file.  The prohibition on code snippets and symbols keeps the file compact and robust to code changes; the one-off task exclusion prevents preference drift from idiosyncratic details.

%......................................................................
\subsection{Pre-Finish Verification}
\label{app:prefinish}
%......................................................................

\begin{prompt}
## Pre-Finish Verification

Before finish(success=True):

1. Re-read and verify every modified file.
1. Run required checks (lint, typecheck, tests); fix failures.
1. Check each user requirement against delivery.
1. If any check fails, keep working.
1. After 3 failed retries of same fix, rethink from scratch.
\end{prompt}

The three-retry threshold forces the model to break out of repetitive loops by abandoning the current approach and reconsidering the problem from first principles.

%......................................................................
\subsection{File Browsing Protocol}
\label{app:file-browsing}
%......................................................................

\begin{prompt}
## File Browsing

Collect info and code snippets in PWD/tmp/file-information-{unique_id}.md without overthinking, then review and think deeply.
\end{prompt}

This instruction applies the same two-phase collect-then-synthesize discipline used for web research (Section~\ref{sec:web-research}) to the local file system.  By writing a structured summary of each file's relevant information into a temporary markdown file, the agent externalizes its understanding into a compact artifact that is typically much smaller than the raw source files, freeing up context window capacity for subsequent reasoning and editing.

%......................................................................
\subsection{Desktop Application Control}
\label{app:desktop}
%......................................................................

\begin{prompt}
## Desktop Apps

Use screenshots, keyboard, and mouse. Don't launch VS Code or its extensions.
\end{prompt}

The prohibition on launching VS~Code prevents a recursive loop: since the agent runs inside a VS~Code extension, launching another instance could corrupt the host session or create deadlocks.

\section{Painless Software Engineering}
\label{app:painless}

A central claim of our system is that natural-language interaction can replace manual code inspection and ad hoc scripting for understanding and evolving nontrivial subsystems.  We illustrate this with a real development session in which the worktree merge workflow (Section~\ref{sec:worktree-agent}) was first understood and then redesigned entirely through conversational prompts.

\paragraph{Step~1: Understanding the existing workflow.} The developer asks the agent to explain the current post-task git lifecycle.  The agent reads the source code and returns a structured summary of the four-phase lifecycle.

\paragraph{Step~2: Simplifying the workflow.} The developer decides the three-way choice (auto-merge, manual merge, discard) is unnecessarily complex and issues a redesign request in natural language.  The agent modifies six files across Python and TypeScript, removes the manual-merge method, simplifies the UI to two buttons, updates tests, and verifies all 28 worktree tests pass.

\paragraph{Step~3: Investigating unexpected state.} After testing, the developer notices files appearing in the Source Control panel and asks why.  The agent traces the execution path and discovers the implementation deliberately unstages changes for user review.

\paragraph{Step~4: Directing a design change.} The developer decides the extra review step is unnecessary and says ``Fix it'' in one sentence.  The agent replaces the unstage call with a conditional commit, adds an edge-case guard, updates one test, and verifies all 104 worktree tests pass.

The developer never opens a source file, never writes a line of code, and never runs a test manually.  This style of development becomes possible because of the agent hierarchy: the Worktree Agent isolates changes on a branch, the Chat Agent preserves context across tasks, and the Relentless Agent continues when the context window is exhausted.

%%%%%%%%%%%%%%%%%%%%%%%%%%%%%%%%%%%%%%%%%%%%%%%%%%%%%%%%%%%%
\section*{NeurIPS Paper Checklist}
%%%%%%%%%%%%%%%%%%%%%%%%%%%%%%%%%%%%%%%%%%%%%%%%%%%%%%%%%%%%

\begin{enumerate}

\item {\bf Claims}
    \item[] Question: Do the main claims made in the abstract and introduction accurately reflect the paper's contributions and scope?
    \item[] Answer: \answerYes{}
    \item[] Justification: The abstract and introduction (Section~\ref{sec:intro}) clearly state the contributions---a five-layer agent architecture, a structured system prompt, a VS~Code extension, and a 62.2\% pass rate on Terminal-Bench~2.0---and the scope is limited to software engineering tasks.
    \item[] Guidelines:
    \begin{itemize}
        \item The answer \answerNA{} means that the abstract and introduction do not include the claims made in the paper.
        \item The abstract and/or introduction should clearly state the claims made, including the contributions made in the paper and important assumptions and limitations. A \answerNo{} or \answerNA{} answer to this question will not be perceived well by the reviewers. 
        \item The claims made should match theoretical and experimental results, and reflect how much the results can be expected to generalize to other settings. 
        \item It is fine to include aspirational goals as motivation as long as it is clear that these goals are not attained by the paper. 
    \end{itemize}

\item {\bf Limitations}
    \item[] Question: Does the paper discuss the limitations of the work performed by the authors?
    \item[] Answer: \answerYes{}
    \item[] Justification: Section~\ref{sec:evaluation} discusses consistently failed tasks and their failure categories (graphical/multimedia requirements, resource-intensive builds, niche domain knowledge). The conclusion acknowledges that results depend on a single frontier model (Claude Opus~4.6) and a single benchmark.
    \item[] Guidelines:
    \begin{itemize}
        \item The answer \answerNA{} means that the paper has no limitation while the answer \answerNo{} means that the paper has limitations, but those are not discussed in the paper. 
        \item The authors are encouraged to create a separate ``Limitations'' section in their paper.
        \item The paper should point out any strong assumptions and how robust the results are to violations of these assumptions (e.g., independence assumptions, noiseless settings, model well-specification, asymptotic approximations only holding locally). The authors should reflect on how these assumptions might be violated in practice and what the implications would be.
        \item The authors should reflect on the scope of the claims made, e.g., if the approach was only tested on a few datasets or with a few runs. In general, empirical results often depend on implicit assumptions, which should be articulated.
        \item The authors should reflect on the factors that influence the performance of the approach. For example, a facial recognition algorithm may perform poorly when image resolution is low or images are taken in low lighting. Or a speech-to-text system might not be used reliably to provide closed captions for online lectures because it fails to handle technical jargon.
        \item The authors should discuss the computational efficiency of the proposed algorithms and how they scale with dataset size.
        \item If applicable, the authors should discuss possible limitations of their approach to address problems of privacy and fairness.
        \item While the authors might fear that complete honesty about limitations might be used by reviewers as grounds for rejection, a worse outcome might be that reviewers discover limitations that aren't acknowledged in the paper. The authors should use their best judgment and recognize that individual actions in favor of transparency play an important role in developing norms that preserve the integrity of the community. Reviewers will be specifically instructed to not penalize honesty concerning limitations.
    \end{itemize}

\item {\bf Theory assumptions and proofs}
    \item[] Question: For each theoretical result, does the paper provide the full set of assumptions and a complete (and correct) proof?
    \item[] Answer: \answerNA{}
    \item[] Justification: The paper does not include theoretical results. It is a systems paper presenting an architecture and empirical evaluation.
    \item[] Guidelines:
    \begin{itemize}
        \item The answer \answerNA{} means that the paper does not include theoretical results. 
        \item All the theorems, formulas, and proofs in the paper should be numbered and cross-referenced.
        \item All assumptions should be clearly stated or referenced in the statement of any theorems.
        \item The proofs can either appear in the main paper or the supplemental material, but if they appear in the supplemental material, the authors are encouraged to provide a short proof sketch to provide intuition. 
        \item Inversely, any informal proof provided in the core of the paper should be complemented by formal proofs provided in appendix or supplemental material.
        \item Theorems and Lemmas that the proof relies upon should be properly referenced. 
    \end{itemize}

    \item {\bf Experimental result reproducibility}
    \item[] Question: Does the paper fully disclose all the information needed to reproduce the main experimental results of the paper to the extent that it affects the main claims and/or conclusions of the paper (regardless of whether the code and data are provided or not)?
    \item[] Answer: \answerYes{}
    \item[] Justification: Section~\ref{sec:evaluation} describes the benchmark (Terminal-Bench~2.0), the evaluation framework (Harbor), the model (Claude Opus~4.6), the hardware (2025 MacBook Air M4, 24\,GB RAM), the number of trials (5 per task), and the 9 skipped tasks. The full agent architecture is described in Section~\ref{sec:architecture} and the system prompt in Section~\ref{sec:system-prompt} and Appendix~\ref{app:full-prompt}.
    \item[] Guidelines:
    \begin{itemize}
        \item The answer \answerNA{} means that the paper does not include experiments.
        \item If the paper includes experiments, a \answerNo{} answer to this question will not be perceived well by the reviewers: Making the paper reproducible is important, regardless of whether the code and data are provided or not.
        \item If the contribution is a dataset and/or model, the authors should describe the steps taken to make their results reproducible or verifiable. 
        \item Depending on the contribution, reproducibility can be accomplished in various ways. For example, if the contribution is a novel architecture, describing the architecture fully might suffice, or if the contribution is a specific model and empirical evaluation, it may be necessary to either make it possible for others to replicate the model with the same dataset, or provide access to the model. In general. releasing code and data is often one good way to accomplish this, but reproducibility can also be provided via detailed instructions for how to replicate the results, access to a hosted model (e.g., in the case of a large language model), releasing of a model checkpoint, or other means that are appropriate to the research performed.
        \item While NeurIPS does not require releasing code, the conference does require all submissions to provide some reasonable avenue for reproducibility, which may depend on the nature of the contribution. For example
        \begin{enumerate}
            \item If the contribution is primarily a new algorithm, the paper should make it clear how to reproduce that algorithm.
            \item If the contribution is primarily a new model architecture, the paper should describe the architecture clearly and fully.
            \item If the contribution is a new model (e.g., a large language model), then there should either be a way to access this model for reproducing the results or a way to reproduce the model (e.g., with an open-source dataset or instructions for how to construct the dataset).
            \item We recognize that reproducibility may be tricky in some cases, in which case authors are welcome to describe the particular way they provide for reproducibility. In the case of closed-source models, it may be that access to the model is limited in some way (e.g., to registered users), but it should be possible for other researchers to have some path to reproducing or verifying the results.
        \end{enumerate}
    \end{itemize}


\item {\bf Open access to data and code}
    \item[] Question: Does the paper provide open access to the data and code, with sufficient instructions to faithfully reproduce the main experimental results, as described in supplemental material?
    \item[] Answer: \answerNo{The code repository link is withheld to preserve double-blind review. The system is open-source and the code will be made available upon acceptance.}
    \item[] Justification: The open-source repository link is withheld to respect double-blind review (stated in Section~\ref{sec:intro}). The system will be publicly released upon acceptance.
    \item[] Guidelines:
    \begin{itemize}
        \item The answer \answerNA{} means that paper does not include experiments requiring code.
        \item Please see the NeurIPS code and data submission guidelines (\url{https://neurips.cc/public/guides/CodeSubmissionPolicy}) for more details.
        \item While we encourage the release of code and data, we understand that this might not be possible, so \answerNo{} is an acceptable answer. Papers cannot be rejected simply for not including code, unless this is central to the contribution (e.g., for a new open-source benchmark).
        \item The instructions should contain the exact command and environment needed to run to reproduce the results. See the NeurIPS code and data submission guidelines (\url{https://neurips.cc/public/guides/CodeSubmissionPolicy}) for more details.
        \item The authors should provide instructions on data access and preparation, including how to access the raw data, preprocessed data, intermediate data, and generated data, etc.
        \item The authors should provide scripts to reproduce all experimental results for the new proposed method and baselines. If only a subset of experiments are reproducible, they should state which ones are omitted from the script and why.
        \item At submission time, to preserve anonymity, the authors should release anonymized versions (if applicable).
        \item Providing as much information as possible in supplemental material (appended to the paper) is recommended, but including URLs to data and code is permitted.
    \end{itemize}


\item {\bf Experimental setting/details}
    \item[] Question: Does the paper specify all the training and test details (e.g., data splits, hyperparameters, how they were chosen, type of optimizer) necessary to understand the results?
    \item[] Answer: \answerYes{}
    \item[] Justification: Section~\ref{sec:evaluation} specifies the benchmark, model, hardware, number of trials, skipped tasks, and evaluation methodology. No training or fine-tuning is performed.
    \item[] Guidelines:
    \begin{itemize}
        \item The answer \answerNA{} means that the paper does not include experiments.
        \item The experimental setting should be presented in the core of the paper to a level of detail that is necessary to appreciate the results and make sense of them.
        \item The full details can be provided either with the code, in appendix, or as supplemental material.
    \end{itemize}

\item {\bf Experiment statistical significance}
    \item[] Question: Does the paper report error bars suitably and correctly defined or other appropriate information about the statistical significance of the experiments?
    \item[] Answer: \answerYes{}
    \item[] Justification: We run 5 independent trials per task (445 total runs) and report pass@any (78.7\%), pass@all (43.8\%), and overall pass rate (62.2\%), as well as the distribution of always-pass, always-fail, and mixed-result tasks. Median and mean cost and duration are reported.
    \item[] Guidelines:
    \begin{itemize}
        \item The answer \answerNA{} means that the paper does not include experiments.
        \item The authors should answer \answerYes{} if the results are accompanied by error bars, confidence intervals, or statistical significance tests, at least for the experiments that support the main claims of the paper.
        \item The factors of variability that the error bars are capturing should be clearly stated (for example, train/test split, initialization, random drawing of some parameter, or overall run with given experimental conditions).
        \item The method for calculating the error bars should be explained (closed form formula, call to a library function, bootstrap, etc.)
        \item The assumptions made should be given (e.g., Normally distributed errors).
        \item It should be clear whether the error bar is the standard deviation or the standard error of the mean.
        \item It is OK to report 1-sigma error bars, but one should state it. The authors should preferably report a 2-sigma error bar than state that they have a 96\% CI, if the hypothesis of Normality of errors is not verified.
        \item For asymmetric distributions, the authors should be careful not to show in tables or figures symmetric error bars that would yield results that are out of range (e.g., negative error rates).
        \item If error bars are reported in tables or plots, the authors should explain in the text how they were calculated and reference the corresponding figures or tables in the text.
    \end{itemize}

\item {\bf Experiments compute resources}
    \item[] Question: For each experiment, does the paper provide sufficient information on the computer resources (type of compute workers, memory, time of execution) needed to reproduce the experiments?
    \item[] Answer: \answerYes{}
    \item[] Justification: Section~\ref{sec:evaluation} specifies the hardware (2025 MacBook Air 15'' with M4 processor and 24\,GB RAM) and reports median/mean cost (\$0.45/\$0.90 per trial) and median/mean duration (202\,s/446\,s per trial).
    \item[] Guidelines:
    \begin{itemize}
        \item The answer \answerNA{} means that the paper does not include experiments.
        \item The paper should indicate the type of compute workers CPU or GPU, internal cluster, or cloud provider, including relevant memory and storage.
        \item The paper should provide the amount of compute required for each of the individual experimental runs as well as estimate the total compute. 
        \item The paper should disclose whether the full research project required more compute than the experiments reported in the paper (e.g., preliminary or failed experiments that didn't make it into the paper). 
    \end{itemize}
    
\item {\bf Code of ethics}
    \item[] Question: Does the research conducted in the paper conform, in every respect, with the NeurIPS Code of Ethics \url{https://neurips.cc/public/EthicsGuidelines}?
    \item[] Answer: \answerYes{}
    \item[] Justification: We have reviewed the NeurIPS Code of Ethics. The research involves building and evaluating a software engineering assistant; no human subjects, sensitive data, or dual-use concerns are involved.
    \item[] Guidelines:
    \begin{itemize}
        \item The answer \answerNA{} means that the authors have not reviewed the NeurIPS Code of Ethics.
        \item If the authors answer \answerNo{}, they should explain the special circumstances that require a deviation from the Code of Ethics.
        \item The authors should make sure to preserve anonymity (e.g., if there is a special consideration due to laws or regulations in their jurisdiction).
    \end{itemize}


\item {\bf Broader impacts}
    \item[] Question: Does the paper discuss both potential positive societal impacts and negative societal impacts of the work performed?
    \item[] Answer: \answerYes{}
    \item[] Justification: The paper discusses positive impacts (making software development more accessible, reducing developer toil). Potential negative impacts include job displacement for routine programming tasks and the risk of generating vulnerable code if the system prompt's verification disciplines are removed. The system mitigates the latter through mandatory pre-finish verification (Appendix~\ref{app:prefinish}) and test-first discipline (Appendix~\ref{app:testing}).
    \item[] Guidelines:
    \begin{itemize}
        \item The answer \answerNA{} means that there is no societal impact of the work performed.
        \item If the authors answer \answerNA{} or \answerNo{}, they should explain why their work has no societal impact or why the paper does not address societal impact.
        \item Examples of negative societal impacts include potential malicious or unintended uses (e.g., disinformation, generating fake profiles, surveillance), fairness considerations (e.g., deployment of technologies that could make decisions that unfairly impact specific groups), privacy considerations, and security considerations.
        \item The conference expects that many papers will be foundational research and not tied to particular applications, let alone deployments. However, if there is a direct path to any negative applications, the authors should point it out. For example, it is legitimate to point out that an improvement in the quality of generative models could be used to generate Deepfakes for disinformation. On the other hand, it is not needed to point out that a generic algorithm for optimizing neural networks could enable people to train models that generate Deepfakes faster.
        \item The authors should consider possible harms that could arise when the technology is being used as intended and functioning correctly, harms that could arise when the technology is being used as intended but gives incorrect results, and harms following from (intentional or unintentional) misuse of the technology.
        \item If there are negative societal impacts, the authors could also discuss possible mitigation strategies (e.g., gated release of models, providing defenses in addition to attacks, mechanisms for monitoring misuse, mechanisms to monitor how a system learns from feedback over time, improving the efficiency and accessibility of ML).
    \end{itemize}
    
\item {\bf Safeguards}
    \item[] Question: Does the paper describe safeguards that have been put in place for responsible release of data or models that have a high risk for misuse (e.g., pre-trained language models, image generators, or scraped datasets)?
    \item[] Answer: \answerNA{}
    \item[] Justification: The paper does not release a pretrained language model or dataset. The system is an agent framework that wraps existing commercial LLM APIs.
    \item[] Guidelines:
    \begin{itemize}
        \item The answer \answerNA{} means that the paper poses no such risks.
        \item Released models that have a high risk for misuse or dual-use should be released with necessary safeguards to allow for controlled use of the model, for example by requiring that users adhere to usage guidelines or restrictions to access the model or implementing safety filters. 
        \item Datasets that have been scraped from the Internet could pose safety risks. The authors should describe how they avoided releasing unsafe images.
        \item We recognize that providing effective safeguards is challenging, and many papers do not require this, but we encourage authors to take this into account and make a best faith effort.
    \end{itemize}

\item {\bf Licenses for existing assets}
    \item[] Question: Are the creators or original owners of assets (e.g., code, data, models), used in the paper, properly credited and are the license and terms of use explicitly mentioned and properly respected?
    \item[] Answer: \answerYes{}
    \item[] Justification: All tools, benchmarks, and models used are cited: Terminal-Bench~2.0, Harbor, Claude Opus~4.6, Playwright, VS~Code, and all related works. The system is open-source.
    \item[] Guidelines:
    \begin{itemize}
        \item The answer \answerNA{} means that the paper does not use existing assets.
        \item The authors should cite the original paper that produced the code package or dataset.
        \item The authors should state which version of the asset is used and, if possible, include a URL.
        \item The name of the license (e.g., CC-BY 4.0) should be included for each asset.
        \item For scraped data from a particular source (e.g., website), the copyright and terms of service of that source should be provided.
        \item If assets are released, the license, copyright information, and terms of use in the package should be provided. For popular datasets, \url{paperswithcode.com/datasets} has curated licenses for some datasets. Their licensing guide can help determine the license of a dataset.
        \item For existing datasets that are re-packaged, both the original license and the license of the derived asset (if it has changed) should be provided.
        \item If this information is not available online, the authors are encouraged to reach out to the asset's creators.
    \end{itemize}

\item {\bf New assets}
    \item[] Question: Are new assets introduced in the paper well documented and is the documentation provided alongside the assets?
    \item[] Answer: \answerNA{}
    \item[] Justification: No new datasets or pretrained models are released. The open-source code will be documented upon release.
    \item[] Guidelines:
    \begin{itemize}
        \item The answer \answerNA{} means that the paper does not release new assets.
        \item Researchers should communicate the details of the dataset/code/model as part of their submissions via structured templates. This includes details about training, license, limitations, etc. 
        \item The paper should discuss whether and how consent was obtained from people whose asset is used.
        \item At submission time, remember to anonymize your assets (if applicable). You can either create an anonymized URL or include an anonymized zip file.
    \end{itemize}

\item {\bf Crowdsourcing and research with human subjects}
    \item[] Question: For crowdsourcing experiments and research with human subjects, does the paper include the full text of instructions given to participants and screenshots, if applicable, as well as details about compensation (if any)? 
    \item[] Answer: \answerNA{}
    \item[] Justification: The paper does not involve crowdsourcing or research with human subjects.
    \item[] Guidelines:
    \begin{itemize}
        \item The answer \answerNA{} means that the paper does not involve crowdsourcing nor research with human subjects.
        \item Including this information in the supplemental material is fine, but if the main contribution of the paper involves human subjects, then as much detail as possible should be included in the main paper. 
        \item According to the NeurIPS Code of Ethics, workers involved in data collection, curation, or other labor should be paid at least the minimum wage in the country of the data collector. 
    \end{itemize}

\item {\bf Institutional review board (IRB) approvals or equivalent for research with human subjects}
    \item[] Question: Does the paper describe potential risks incurred by study participants, whether such risks were disclosed to the subjects, and whether Institutional Review Board (IRB) approvals (or an equivalent approval/review based on the requirements of your country or institution) were obtained?
    \item[] Answer: \answerNA{}
    \item[] Justification: The paper does not involve human subjects research.
    \item[] Guidelines:
    \begin{itemize}
        \item The answer \answerNA{} means that the paper does not involve crowdsourcing nor research with human subjects.
        \item Depending on the country in which research is conducted, IRB approval (or equivalent) may be required for any human subjects research. If you obtained IRB approval, you should clearly state this in the paper. 
        \item We recognize that the procedures for this may vary significantly between institutions and locations, and we expect authors to adhere to the NeurIPS Code of Ethics and the guidelines for their institution. 
        \item For initial submissions, do not include any information that would break anonymity (if applicable), such as the institution conducting the review.
    \end{itemize}

\item {\bf Declaration of LLM usage}
    \item[] Question: Does the paper describe the usage of LLMs if it is an important, original, or non-standard component of the core methods in this research? Note that if the LLM is used only for writing, editing, or formatting purposes and does \emph{not} impact the core methodology, scientific rigor, or originality of the research, declaration is not required.
    \item[] Answer: \answerYes{}
    \item[] Justification: The entire system is built around LLM usage. Section~\ref{sec:architecture} describes how LLMs are used as the core reasoning engine via native function calling (Section~\ref{sec:ks-agent}). The paper also discloses that the system was built using itself (Section~\ref{sec:intro}).
    \item[] Guidelines:
    \begin{itemize}
        \item The answer \answerNA{} means that the core method development in this research does not involve LLMs as any important, original, or non-standard components.
        \item Please refer to our LLM policy in the NeurIPS handbook for what should or should not be described.
    \end{itemize}

\end{enumerate}

\end{document}
